\section{The finite hypergeometric algebra}\label{sec:series}

\subsection{Rising factorials and coefficients}

\begin{definition}[Rising factorial]\label{def:pochhammer}
For an element \(x\) of a commutative ring, define
\[
  \poch{x}{0}=1,
  \qquad
  \poch{x}{n+1}=\poch{x}{n}(x+n).
\]
For coded rationals or complex rationals this is a primitive recursive finite
product.
\end{definition}

\begin{lemma}[Elementary Pochhammer identities]\label{lem:pochhammer}
Whenever the displayed operations are defined,
\begin{align}
 \poch{x}{m+n}&=\poch{x}{m}\poch{x+m}{n},\label{eq:poch-split}\\
 \poch{x}{n+1}&=x\poch{x+1}{n},\label{eq:poch-shift}\\
 \poch{-m}{n}&=(-1)^n\frac{m!}{(m-n)!}\quad(n\le m),\label{eq:poch-neg}\\
 \poch{-m}{n}&=0\quad(n>m).\label{eq:poch-zero}
\end{align}
All four are provable in iPRA by quantifier-free induction.
\end{lemma}

\begin{proof}
For \eqref{eq:poch-split}, induct on \(n\).  The successor step multiplies
both sides by \(x+m+n\).  Formula \eqref{eq:poch-shift} is the case in which
the first factor is separated.  For \eqref{eq:poch-neg}, write each factor
\(-m+j=-(m-j)\) and induct while maintaining \(n\le m\); cross-multiplication
avoids division.  At \(n=m+1\), the factor \(-m+m\) is zero, and every later
product remains zero.  These are polynomial identities in integer codes.
\end{proof}

\begin{definition}[Gauss coefficient]\label{def:gauss-coeff}
Assume \(\poch{c}{n}\ne0\).  Define
\[
 A_n(a,b;c)=\frac{\poch{a}{n}\poch{b}{n}}
                 {\poch{c}{n}\,n!}.
\]
The guarded total definition uses \(\operatorname{qdiv}\).  The partial sum is
\[
 H_N(a,b;c;z)=\sum_{n<N}A_n(a,b;c)z^n.
\]
\end{definition}

\begin{lemma}[Coefficient ratio]\label{lem:ratio}
Under the denominator guards,
\begin{equation}\label{eq:ratio}
 (n+1)(n+c)A_{n+1}=(n+a)(n+b)A_n.
\end{equation}
If \((n+1)(n+c)\) is apart from zero, this is equivalently
\[
  \frac{A_{n+1}}{A_n}=\frac{(n+a)(n+b)}{(n+1)(n+c)}
\]
when \(A_n\ne0\).
\end{lemma}

\begin{proof}
Expand the four successor recurrences for the rising factorials and factorial,
then cancel only after multiplying through by the denominators.  The
cross-multiplied identity \eqref{eq:ratio} is valid even when \(A_n=0\).
\end{proof}

\begin{definition}[Formal Gauss series]\label{def:formal-gauss}
The formal object
\[
   \F{a}{b}{c}{z}:=\sum_{n\ge0}A_n(a,b;c)z^n
\]
means the primitive recursive coefficient table \(n\mapsto A_n\), together
with the guards \(\poch c n\ne0\) needed at the observed indices.
\end{definition}

\begin{proposition}[Symmetry and normalization]\label{prop:symmetry}
Formally,
\[
 \F{a}{b}{c}{z}=\F{b}{a}{c}{z},
 \qquad
 \F{a}{b}{c}{0}=1.
\]
\end{proposition}

\begin{proof}
The coefficient formula is symmetric in \(a,b\), and the constant coefficient
is \(1\).  No limit is involved.
\end{proof}

\subsection{Termination and the polynomial formula}

\begin{theorem}[Terminating series]\label{thm:termination}
For \(m\in\N\),
\begin{equation}\label{eq:terminating}
 \F{-m}{b}{c}{z}
   =\sum_{n=0}^{m}(-1)^n\binom mn\frac{\poch b n}{\poch c n}z^n
\end{equation}
as a polynomial identity, provided the finitely many denominator factors are
nonzero.
\end{theorem}

\begin{proof}
By \cref{lem:pochhammer},
\[
 \frac{\poch{-m}{n}}{n!}
 =(-1)^n\frac{m!}{n!(m-n)!}=(-1)^n\binom mn
 \quad(n\le m),
\]
and every coefficient with \(n>m\) vanishes.  Cross-multiplying by the finite
product \(\prod_{j<m}(c+j)\) turns the whole statement into an integer
polynomial identity.  The proof uses bounded products and induction only.
\end{proof}

\begin{corollary}[Binomial specialization]\label{cor:binomial-termination}
For \(m\in\N\),
\[
 \F{-m}{b}{b}{z}=(1-z)^m.
\]
In particular, \((1+z)^m=\F{-m}{b}{b}{-z}\).
\end{corollary}

\begin{proof}
The factors \(\poch b n\) cancel under the supplied guards, and
\eqref{eq:terminating} becomes the finite binomial theorem.  The binomial
theorem itself is a quantifier-free induction on \(m\).
\end{proof}


\subsection{Regularization at a singular denominator parameter}

The ordinary series is singular when \(c\) approaches a nonpositive integer,
but multiplication by the reciprocal gamma function removes the singularity.

\begin{theorem}[Regularized limit at \(c=-m\)]\label{thm:regularized-c}
Let \(m\in\N\).  Coefficientwise, and uniformly on every strict disk after
supplying the effective reciprocal-gamma bounds,
\begin{equation}\label{eq:regularized-c}
 \lim_{c\to-m}\frac{\F a b c z}{\Gamma(c)}
 =\frac{\poch a{m+1}\poch b{m+1}}{(m+1)!}z^{m+1}
  \F{a+m+1}{b+m+1}{m+2}{z}.
\end{equation}
\end{theorem}

\begin{proof}
Use the exact gamma recurrence identity
\[
 \frac1{\Gamma(c)\poch c n}=\frac1{\Gamma(c+n)}.
\]
Thus the coefficient of \(z^n\) on the left is
\[
 \frac{\poch a n\poch b n}{n!\,\Gamma(c+n)}.
\]
The reciprocal gamma function has a zero at every nonpositive integer.  At
\(c=-m\), the coefficients with \(n\le m\) therefore tend to zero.  Write
\(n=m+1+k\) for the remaining coefficients.  Since
\(\Gamma(k+1)=k!\), their limits are
\[
 \frac{\poch a{m+1+k}\poch b{m+1+k}}
      {(m+1+k)!\,k!}.
\]
Splitting both rising factorials at \(m+1\) and using
\((m+2)_k=(m+1+k)!/(m+1)!\) gives exactly the coefficient on the right of
\eqref{eq:regularized-c}.  This is a finite identity for each coefficient.
For uniform convergence on \(\norm z\le r<1\), choose a small rational
parameter disk around \(-m\), a common ratio cutoff for its tail, and then a
finite head on which the reciprocal-gamma name converges effectively.  The
gamma packet supplies the simple-zero and local modulus data.
\end{proof}

\subsection{Generalized hypergeometric coefficients}

\begin{definition}[Generalized coefficient]\label{def:pfq}
For finite parameter vectors
\(\boldsymbol a=(a_1,\ldots,a_p)\) and
\(\boldsymbol b=(b_1,\ldots,b_q)\), put
\[
 T_n(\boldsymbol a;\boldsymbol b;z)
  =\frac{\prod_{i<p}\poch{a_i}{n}}
         {n!\prod_{j<q}\poch{b_j}{n}}z^n.
\]
The associated formal series is \({}_pF_q\).  All vector products are bounded
folds over the externally fixed or coded finite lengths \(p,q\).
\end{definition}

\begin{proposition}[Hypergeometric term ratio]\label{prop:pfq-ratio}
The successive terms satisfy the cross-multiplied identity
\[
 (n+1)\prod_{j<q}(n+b_j)T_{n+1}
 =z\prod_{i<p}(n+a_i)T_n.
\]
Thus the ratio of successive nonzero terms is a rational function of \(n\),
as asserted on the MathWorld page.
\end{proposition}

\begin{proof}
Apply the successor equation for every Pochhammer factor.  Bounded-product
congruence and associativity suffice.
\end{proof}

\subsection{Formal differentiation}

For a formal series \(F(z)=\sum_{n\ge0}f_nz^n\), define
\(DF=\sum_{n\ge0}(n+1)f_{n+1}z^n\).  This is a primitive recursive operation
on coefficient tables.

\begin{theorem}[Derivative formulas]\label{thm:derivatives}
For every \(r\in\N\) for which the denominator guards hold,
\begin{equation}\label{eq:nth-derivative}
 D^r\F{a}{b}{c}{z}
 =\frac{\poch a r\poch b r}{\poch c r}
   \F{a+r}{b+r}{c+r}{z}.
\end{equation}
The theorem is uniform in \(r\) when the shifted coefficient table and guards
are jointly primitive recursive.
\end{theorem}

\begin{proof}
For \(r=1\), the coefficient of \(z^n\) on the left is
\[
 (n+1)\frac{\poch a{n+1}\poch b{n+1}}
                   {\poch c{n+1}(n+1)!}
 =\frac{ab}{c}\frac{\poch{a+1}n\poch{b+1}n}
                         {\poch{c+1}n n!},
\]
using \eqref{eq:poch-shift}.  This is an exact rational identity.  Induction
on \(r\), together with
\(\poch a{r+1}=\poch a r(a+r)\), gives
\eqref{eq:nth-derivative}.  The induction invariant is a single
cross-multiplied coefficient equality.
\end{proof}

\subsection{The coefficient recurrence as a characterization}

\begin{theorem}[Formal uniqueness]\label{thm:formal-uniqueness}
Let \((B_n)\) be a coefficient table satisfying
\[
 B_0=1,
 \qquad
 (n+1)(n+c)B_{n+1}=(n+a)(n+b)B_n,
\]
and suppose each \((n+1)(n+c)\) up to the observed index is apart from zero.
Then \(B_n=A_n(a,b;c)\) for every \(n\).
\end{theorem}

\begin{proof}
Induct on \(n\).  The base is \(1\).  In the successor step, both \(A_{n+1}\)
and \(B_{n+1}\) solve the same linear equation with an apart coefficient.
Cross-multiply and use cancellation in the coded rational field.  In the
object theory one may avoid field cancellation by carrying numerator and
denominator representatives and the nonzero product as a side assumption.
\end{proof}

\begin{remark}[Why uniqueness is powerful]
To prove a transformation, it often suffices to show that a transformed
series has constant term \(1\) and satisfies the same recurrence.  This turns
an identity between infinite objects into a finite calculation with the
chain rule, followed by \cref{thm:formal-uniqueness}.
\end{remark}

\subsection{Finite verification principle}

\begin{proposition}[Certificate checker for identities]\label{prop:certificate-checker}
Fix rational parameter maps
\((a,b,c,z)\mapsto(a',b',c',\phi(z))\) and a prefactor \(P(z)\) whose formal
coefficients are primitive recursive.  To prove
\[
 \F a b c z=P(z)\F{a'}{b'}{c'}{\phi(z)}
\]
it is enough to provide either:
\begin{enumerate}
\item a coefficient identity for the \(n\)-th coefficient, expressed as a
finite sum and proved by induction in \(n\); or
\item an algebraic certificate that the right side satisfies the same
second-order differential equation and the same first two initial
coefficients as the left side.
\end{enumerate}
Both certificate formats have primitive recursive checkers.
\end{proposition}

\begin{proof}
In the first format, the checker expands products and substitutions in
\(R_{n+1}\).  In the second, the differential expression has coefficients
obtained by a finite number of Cauchy products and derivatives; equality to
zero at each coefficient is decidable.  The ODE recurrence then determines
all higher coefficients from the first two, or from the constant coefficient
for the normalized regular solution.  The checker never manipulates an
actual completed infinite series.
\end{proof}
